\section{Additional Related Work}
\label{app:related_work}

\paragraph{Empirical studies: probing and editing LLM knowledge.}
\citet{geva2021transformerfeedforwardlayerskeyvalue,geva2022transformerfeedforwardlayersbuild} observed that knowledge is often stored within MLPs via key--value mappings, motivating a line of work that attempts to reverse engineer the facts encoded in MLPs~\citep{dai2022knowledgeneuronspretrainedtransformers,nanda2023factfinding} and to edit them~\citep{dai2022knowledgeneuronspretrainedtransformers,meng2023locatingeditingfactualassociations,memit,model_edit_scaling,gu2024modeleditingharmsgeneral,fang2025alphaeditnullspaceconstrainedknowledge,sun2025mitigatingnegativeinterferencemultilingual}. These studies provide strong empirical evidence that MLPs act as a locus of factual storage in large language models.

\paragraph{Empirical studies: scaling factual knowledge.}
A related empirical line of work formalizes factual knowledge as associative recall over key--value stores and studies its scaling behavior~\citep{elhage2022toymodelssuperposition,allenzhu2024physicslanguagemodels33,zucchet2025languagemodelslearnfacts}. These works consistently find that trained models store facts at the asymptotically optimal rate implied by \Cref{thm: info_bounds_capacity-const}~\citep{allenzhu2024physicslanguagemodels33,zucchet2025languagemodelslearnfacts,morris2025languagemodelsmemorize}, which motivates the search for explicit constructions with comparable parameter efficiency.
